\subsection{A cohomological obstruction for the numerical tails}

The coefficient deformation in A-015 has natural ideals with exactly
the required associated graded modules. We now determine all their
filtered bilinear pairings. This tests a proposed bridge already on
cohomology, so allowing higher homotopies in a cochain multiplication
does not remove this particular obstruction.

\begin{theorem}[A-023: \status{Proved}; all filtered pairings of numerical tails]
\label{sa08:A023}
Put $R=\operatorname{Int}(\Z)\subset\Q[T]$, with polynomial-degree
filtration, and write $\beta_a=\binom{T}{a}$. For $N\geq0$ let
\[
 J_N=\{f\in R:f(0)=\cdots=f(N-1)=0\}
     =\bigoplus_{a\geq N}\Z\beta_a.
\]
An empty list of vanishing conditions gives $J_0=R$.
Then $\operatorname{gr}J_N=I_N$ as graded modules over
$\operatorname{gr}R=\Gamma_{\Z}(T)$, but for every $N,M,L\geq0$
the full group of $R$-bilinear pairings
\[
 b:J_N\times J_M\longrightarrow J_L,
 \qquad b(F_aJ_N,F_bJ_M)\subset F_{a+b}J_L,
\]
is
\[
 \begin{cases}
  \Z\cdot[(f,g)\mapsto fg],&L\leq\max(N,M),\\
  0,&L>\max(N,M).
 \end{cases}
\]
Consequently, for $N,M>0$ no such pairing
$J_N\times J_M\to J_{N+M}$ lifts the nonzero PD multiplication
$I_N\otimes I_M\to I_{N+M}$.
In particular, no filtered cohomology algebra with coefficient ring
$R$ and specified filtered coefficient modules
$H^i(X;\Z)\otimes J_i$ can have the PD-Rees associated graded
algebra for every finite simplicial set $X$. This conclusion applies
to cohomology algebras of homotopy-coherent models as well, whenever
these coefficient-module and associated-graded identifications are
part of the proposed comparison.
\end{theorem}

\begin{proof}
The binomial basis and its unitriangular evaluations at the
nonnegative integers show inductively that the first $N$ binomial
coefficients of $f$ vanish exactly when $f(0),\ldots,f(N-1)$ do.
Thus $J_N$ has the displayed basis. It is an ideal because its
definition is by evaluation kernels. Its leading symbols identify
its associated graded with $I_N$; the module action is the binomial
multiplication of A-015.

The fraction field of $R$ is $K=\Q(T)$. Each nonzero ideal $J_N$
becomes the one-dimensional $K$-space $K$ after scalar extension.
Therefore an $R$-bilinear $b$ becomes multiplication by one rational
function $q(T)\in K$:
\[
                    b(f,g)=q(T)f(T)g(T).
\]
This equality holds before localization because the target $J_L$
is a submodule of the field $K$; any torsion in the source tensor
product is necessarily killed. Applying the filtration bound to
$\beta_N,\beta_M$ shows that $q=0$ or
$\deg_\infty q\leq0$.

For every integer $n\geq\max(N,M)$ beyond the finitely many poles,
both inputs $\beta_n$ lie in the required ideals and
$\beta_n(n)=1$. Hence
\[
              q(n)=b(\beta_n,\beta_n)(n)\in\Z.
\]
A rational function of degree at infinity at most zero has a finite
limit along $n\to+\infty$. A convergent sequence of integers is
eventually constant. It follows that $q(n)=c$ for infinitely many
$n$, for some $c\in\Z$, and therefore $q(T)=c$ identically.
This proves that every pairing is ordinary multiplication by an
integer. If $L\leq\max(N,M)$, one input vanishes at every
$0\leq j<L$, so all such integer multiples do take values in $J_L$.
If $L>\max(N,M)$, set $l=\max(N,M)$ and use
$\beta_l(l)=1$ in both inputs. The output vanishes at $l$, forcing
$c=0$. This proves the classification.

For the last assertion it suffices to use the finite simplicial
torus $X=(\Delta[1]/\partial\Delta[1])^2$. It has one vertex,
three nondegenerate edges $a,b,c$, and two nondegenerate triangles
$u,v$, with
\[
 \partial u=\partial v=a+b-c.
\]
The cocycles $\alpha=(1,0,1)$ and $\beta=(0,1,1)$ in this edge
basis have Alexander--Whitney product with triangle values $(1,0)$.
Its value on the fundamental cycle $u-v$ is one. Thus the cross
product between the two copies of $H^1(X;\Z)$ and $H^2(X;\Z)$
would induce an $R$-bilinear filtered pairing
$J_1\times J_1\to J_2$ whose associated graded is PD
multiplication. The classification says it is zero, whereas that
associated graded sends $(T^{[1]},T^{[1]})$ to $2T^{[2]}\ne0$.

There is also a short finite-weight certificate for this normalized
pairing. Since $F_1J_2=0$, its required leading term forces
$b(T,T)=2\beta_2=T(T-1)$. Bilinearity and
$2\beta_2=T(T-1)$ give
\[
       4b(\beta_2,\beta_2)=T(T-1)^3.
\]
Evaluating at $T=2$ would give $4b(\beta_2,\beta_2)(2)=2$, which
is impossible in $\Z$ and already contradictory modulo four.
The program \texttt{A\_numerical\_tail\_pairings.py} verifies the
complete finite torus chain and cup-product data, this polynomial
identity, and all four possible residues in the final contradiction.
\end{proof}

This is a theorem about the specified numerical tails. A filtered
model whose cohomology deforms $I_i$ by different modules, or whose
filtration has nonzero spectral-sequence differentials, is not
excluded by this statement alone. Statement~\ref{sa08:A026}
removes the prescribed-lattice hypothesis and proves that naturality
and the coefficient maps force the relevant circle and torus
filtrations to have no such differentials.
